\noindent\textbf{Elección del procedimiento.} Definir $S_g=Mg\sin\varphi$ y $k_D=\rho C_DA/2$. A potencia constante, subida y bajada cumplen $P=S_gV_1+k_DV_1^3$ y $P=-S_gV_2+k_DV_2^3$. Restar elimina $P$ y da $k_D=S_g(V_1+V_2)/(V_2^3-V_1^3)$.

\begin{enumerate}

\item Convertir pendiente del 3 por ciento en seno.

\[\fitmath{\mathrm{Sg}=\frac{M\cdot g\cdot \mathrm{slope}}{\sqrt{1+{\mathrm{slope}}^{2}}}}\]

\[\fitmath{\mathrm{Sg}=\frac{(1600)\cdot (9.81)\cdot (0.03)}{\sqrt{1+{(0.03)}^{2}}}\approx 470.6682\,\mathrm{N}}\]

\item Eliminar la potencia común entre los dos equilibrios.

\[\fitmath{\mathrm{kD}=\frac{\mathrm{Sg}\cdot \left(V_{1}+V_{2}\right)}{{V_{2}}^{3}-{V_{1}}^{3}}}\]

\[\fitmath{\mathrm{kD}=\frac{(470.6682469)\cdot \left((32.8)+(50.7)\right)}{{(50.7)}^{3}-{(32.8)}^{3}}\approx 0.4135346\,\mathrm{kg/m}}\]

\item Pasar del coeficiente de fuerza cuadrática al aerodinámico.

\[\fitmath{C_D=\frac{2\cdot \mathrm{kD}}{\rho\cdot A}}\]

\[\fitmath{C_D=\frac{2\cdot (0.4135346424)}{(1.2)\cdot (2.3)}\approx 0.2996628\,}\]

\item Recuperar la potencia sustituyendo en el equilibrio de subida.

\[\fitmath{P=\mathrm{Sg}\cdot V_{1}+\mathrm{kD}\cdot {V_{1}}^{3}}\]

\[\fitmath{P=(470.6682469)\cdot (32.8)+(0.4135346424)\cdot {(32.8)}^{3}\approx 30030.54\,\mathrm{W}}\]

\end{enumerate}
